\uuid{X9wY}
\titre{Système linéaire à paramètres/Méthode du Pivot de Gauss}
\auteur{hanani ; mzari}
\organisation{exo7}
\licence{CC BY-NC-SA 4.0}
\datecreate{2020-01-20}
\niveau{L1}
\module{Algèbre}
\chapitre{Déterminant, système linéaire}
\sousChapitre{Système linéaire, rang}
\difficulte{3}
\theme{Systèmes d'équations linéaires, Système linéaire à paramètres, Méthode du Pivot de Gauss}
\exo7qcmid{lille-2.systemes.016}

\contenu{
\begin{question}
\qclassification{200.03}
\qtags{motcle=Système linéaire à paramètres/Méthode du Pivot de Gauss}
On considère le système d'équations, d'inconnue $(x,y,z)\in \Rr^3$ et de paramètres des réels  $a,b$ et $c$ :
$$(\mathtt{S})
\left\{\begin{array}{rcc}
x+ay+a^2z&=&0\\
x+by+b^2z&=&0\\
x+cy+c^2z&=&0.\end{array}\right.$$
Quelles sont les assertions vraies ?
\begin{answers}
\good{$(\mathtt{S}) \Leftrightarrow
\left\{\begin{array}{rcc}
x+ay+a^2z&=&0\\
(b-a)y+(b^2-a^2)z&=&0\\
(c-a)y+(c^2-a^2)z&=&0.\\
\end{array}\right.$}
\bad{Si $a,b$ et $c$ sont des réels deux à deux distincts,  $(\mathtt{S})$ admet une infinité de solutions.}
\bad{Si $a=b$ et $a\neq c$,  $(\mathtt{S})$ admet une unique  solution.}
\bad{$b=c$ et $a\neq c$, $(\mathtt{S})$ n'admet pas de solution.}
\end{answers}
\begin{explanations}
$$(\mathtt{S}) \Leftrightarrow
\left\{\begin{array}{rcc}
x+ay+a^2z&=&0\\
(b-a)y+(b^2-a^2)z&=&0\\
(c-a)y+(c^2-a^2)z&=&0.\\
\end{array}\right.$$
\begin{enumerate}
\item[-]Si $a=b=c$, $(\mathtt{S})\Leftrightarrow  x+ay+a^2z=0 $, $(\mathtt{S})$ admet donc une infinité de solutions.
\item[-]Si $a=b$ et $a\neq c$, ou $a=c$ et $a\neq b$, ou  $b=c$ et $b\neq a$, $(\mathtt{S})$ admet  une infinité de solutions.
\item[-]Si $a,b,c$ sont deux à deux distincts, $(\mathtt{S})$ admet donc une unique solution : $(0,0,0)$.
\end{enumerate}
\end{explanations}
\end{question}
}
